\chapter{The command line}
\label{ch:cli}

BAG is one program, \cmd{bag}, with a command for each step. \cmd{hdl-bag} is the same
program under another name. \cmd{bag --version} prints the version and
\cmd{bag <command> --help} the options of a command.

A \term{spec} is a workbook (\file{.xlsx}) or an IR (\file{.ir.yaml} or \file{.yml}).
An \term{output directory} is a folder that \cmd{bag generate} wrote; the commands that
run EDA tools read its \file{bag_outputs.json} to find the files.

\section{Commands}

\begin{entry}{\cmd{bag template} [\opt{--sample}] \opt{-o} WORKBOOK}
\item[Does] Writes a blank spec workbook to WORKBOOK. With \opt{--sample}, writes the
  sample design \code{demo_ss} instead, and the file header it names
  (\file{header.txt}) in the same folder. An existing \file{header.txt} is kept.
\item[Options] \opt{--output} (\opt{-o}): the workbook to write; required.
  \opt{--sample}: write the sample.
\end{entry}

\begin{entry}{\cmd{bag check} WORKBOOK}
\item[Does] Reads WORKBOOK and runs every check (chapter~\ref{ch:checks}). Prints each
  finding and a summary line. File header paths are relative to the workbook's folder.
\item[Exit status] 1 on any error; warnings alone give 0.
\end{entry}

\begin{entry}{\cmd{bag ir} WORKBOOK \opt{-o} IR}
\item[Does] Checks WORKBOOK as \cmd{bag check} does, then writes its IR
  (chapter~\ref{ch:ir}).
\item[Options] \opt{--output} (\opt{-o}): the IR file to write; required.
\end{entry}

\begin{entry}{\cmd{bag generate} SPEC [\opt{-o} DIR] [\opt{--style} STYLE]}
\item[Does] Checks SPEC, then writes the RTL, the wrapper and the collateral into the
  output directory (chapter~\ref{ch:outputs}). Prints the number of files and the top
  module. An IR is checked as the workbook it converts to.
\item[Options] \opt{--output} (\opt{-o}): the output directory. Without it, BAG uses
  Project > \setting{Output directory}, relative to the folder of SPEC.
  \opt{--style}: \code{flat} or \code{sv-interface}, in place of Wrapper >
  \setting{Port style}.
\end{entry}

\begin{entry}{\cmd{bag lint} DIR}
\item[Does] Lints the generated RTL in DIR with the tools of Project >
  \setting{Lint tools} (section~\ref{sec:lint}).
\item[Exit status] 1 on any finding that is not waived.
\end{entry}

\begin{entry}{\cmd{bag sim} DIR [\opt{--filter} REGEX] [\opt{--ready-shim}]}
\item[Does] Runs the generated cocotb tests in DIR with Verilator
  (section~\ref{sec:sim}). Needs the \code{sim} extra.
\item[Options] \opt{--filter}: run only the tests whose names this regular expression
  finds. \opt{--ready-shim}: run the traffic through a top that drives each AXI READY
  from its VALID; it cannot be used with \opt{--filter}.
\item[Exit status] 1 on any failed test, or when DIR has no generated tests.
\end{entry}

\begin{entry}{\cmd{bag formal} DIR [\opt{--timeout} SECONDS]}
\item[Does] Runs every generated SymbiYosys job in DIR, one after the other
  (section~\ref{sec:formal}).
\item[Options] \opt{--timeout}: the seconds allowed per job, at least 1. A job without a
  verdict in that time is a finding.
\item[Exit status] 1 on any finding: a failed assertion, an unreached cover, a tool
  error or a timeout.
\end{entry}

\begin{entry}{\cmd{bag flow} SPEC [\opt{-o} DIR] [\opt{--style} STYLE] [\opt{--formal} [\opt{--formal-timeout} SECONDS]]}
\item[Does] Runs the steps in turn: check, generate, lint, compile, the generated tests,
  and the formal jobs with \opt{--formal}. Stops at the first step that fails. Formal is
  left out unless you ask for it, because the proofs of a large design take long; run
  \cmd{bag formal} on the output directory later instead.
\item[Options] \opt{--output} (\opt{-o}) and \opt{--style}: as for \cmd{bag generate}.
  \opt{--formal}: also run the formal jobs. \opt{--formal-timeout}: the seconds allowed
  per formal job; it needs \opt{--formal}.
\item[Without tests] For a design with no generated tests (Project >
  \setting{Generate cocotb unit tests} is N, or a user-supplied fabric), the compile
  step also builds the model and the simulation is left out, with a note.
\end{entry}

\section{Exit status}

\begin{description}
\item[0] Everything passed. Warnings and notes do not change this.
\item[1] A check gave an error, a step reported a finding, or a step could not run (for
  example a missing tool, or an output directory without generated tests).
\item[2] The command line itself is wrong: an unknown option, a missing argument, or a
  file that does not exist. The message starts with \code{Usage:}.
\end{description}
